\begin{figure}[H]
\centering

\begin{tikzpicture}[
    scale=0.8,
    transform shape,
    >=stealth,
    neuron/.style={circle, minimum size=9mm, inner sep=0pt},
    hnode/.style={circle, minimum size=9mm, inner sep=0pt},
    outnode/.style={
        circle,
        minimum width=13mm,
        minimum height=13mm,
        inner sep=0pt,
        fill=red!20
    },
    darkarrow/.style={->, draw=darkgrey, line width=0.6pt}
]

\definecolor{darkgrey}{RGB}{90,90,90}

% =====================================
% Latent layer (left)
% =====================================
\node[neuron, fill=red!70] (tau) at (0, 6.0) {$\tau$};

\node[neuron, fill=red!70] (z1) at (0, 4.2) {$z_1$};
\node                      (zvd) at (0, 3.0) {$\vdots$};
\node[neuron, fill=red!70] (zl) at (0, 1.8) {$z_\ell$};

% =====================================
% X positions
% =====================================
\def\xHidden{4.3}
\def\xOut{8.0}

% =====================================
% Y positions of stacked networks
% =====================================
\def\yG{5.0}
\def\yK{2.5}
\def\yH{0.5}
\def\yR{-2.5}

% =====================================
% Hidden layers
% =====================================
\def\HiddenSep{0.9}

\newcommand{\HiddenRange}[4]{%
  \node[hnode, fill=red!40] (#1hTop)  at (\xHidden, {#2+\HiddenSep}) {#3};
  \node                      (#1hDots) at (\xHidden, {#2}) {$\vdots$};
  \node[hnode, fill=red!40] (#1hBot)  at (\xHidden, {#2-\HiddenSep}) {#4};
}

% G: x_{1,G}^(1) ... x_{1,G}^(10)
\HiddenRange{G}{\yG}{$x_{1,\mathcal{G}}^{(1)}$}{$x_{1,\mathcal{G}}^{(10)}$}
% H: x_{1,H}^(1) ... x_{1,H}^(4)
\HiddenRange{H}{\yH}{$x_{1,\mathcal{H}}^{(1)}$}{$x_{1,\mathcal{H}}^{(4)}$}
% R: x_{1,R}^(1) ... x_{1,R}^(4)
\HiddenRange{R}{\yR}{$x_{1,\mathcal{R}}^{(1)}$}{$x_{1,\mathcal{R}}^{(4)}$}

% =====================================
% Output nodes
% =====================================
\node[outnode] (Gout) at (\xOut, \yG) {$\mathcal{G}(z,\tau)$};
\node[outnode] (Kout) at (\xOut, \yK) {$\mathcal{K}(z)$};
\node[outnode] (Hout) at (\xOut, \yH) {$\mathcal{H}(z)$};
\node[outnode] (Rout) at (\xOut, \yR) {$\mathcal{R}(z)$};

% =====================================
% Latent -> G hidden
% =====================================
\draw[darkarrow] (z1.east) to[out=20, in=180] (GhTop.north west);
\draw[darkarrow] (zl.east) to[out=35, in=180] (GhTop.south west);

\draw[darkarrow] (z1.east) to[out=5,  in=180] (GhBot.north west);
\draw[darkarrow] (zl.east) to[out=15, in=180] (GhBot.south west);

% tau -> G hidden
\draw[darkarrow] (tau.east) to[out=-10, in=170] (GhTop.west);
\draw[darkarrow] (tau.east) to[out=-25, in=170] (GhBot.west);

% =====================================
% Latent -> K output (direct, no hidden layer)
% =====================================
\draw[darkarrow] (z1.east) to[out=-20, in=180] (Kout.west);
\draw[darkarrow] (zl.east) to[out=-20, in=180] (Kout.west);

% =====================================
% Latent -> H hidden
% =====================================
\draw[darkarrow] (z1.east) to[out=-25, in=180] (HhTop.north west);
\draw[darkarrow] (zl.east) to[out=-10, in=180] (HhTop.south west);
\draw[darkarrow] (z1.east) to[out=-40, in=180] (HhBot.north west);
\draw[darkarrow] (zl.east) to[out=-25, in=180] (HhBot.south west);

% =====================================
% Latent -> R hidden
% =====================================
\draw[darkarrow] (z1.east) to[out=-45, in=180] (RhTop.north west);
\draw[darkarrow] (zl.east) to[out=-35, in=180] (RhTop.south west);

\draw[darkarrow] (z1.east) to[out=-60, in=180] (RhBot.north west);
\draw[darkarrow] (zl.east) to[out=-50, in=180] (RhBot.south west);

% =====================================
% Hidden -> outputs
% =====================================
\draw[darkarrow] (GhTop.east) -- (Gout.west);
\draw[darkarrow] (GhBot.east) -- (Gout.west);

\draw[darkarrow] (HhTop.east) -- (Hout.west);
\draw[darkarrow] (HhBot.east) -- (Hout.west);

\draw[darkarrow] (RhTop.east) -- (Rout.west);
\draw[darkarrow] (RhBot.east) -- (Rout.west);

\end{tikzpicture}
\caption{The latent factors $z = (z_1, \ldots, z_\ell)^\top$ feed into four neural networks with different architectures and roles. $\mathcal{K}(z)$ is a linear map with no hidden layer. $\mathcal{H}(z)$ and $\mathcal{R}(z)$ each pass through a hidden layer of width 4 to produce their respective outputs. $\mathcal{G}(z, \tau)$ additionally takes the maturity $\tau$ as input and passes through a hidden layer of width 10.}
\label{fig:rolf_architecture_details}

\end{figure}